% =====================================================================
%  31.4  Which edges can be in an MST -- classification and uniqueness.
%  Exports: sec:c17-classify  thm:c17-intersect  thm:c17-cycle  def:c17-forced
%  lem:c17-group  lst:c17-group  prob:c17-unify  fig:c17-unify
%  lst:c17-unify  fig:c17-groups
%  Solved CF 1108F "MST Unification" (2100, L5), live-verified against the
%  CF problem page (title "F. MST Unification", 3s TL, 256 MB, n,m<=2e5,
%  w<=1e9).  Answer = sum over equal-weight Kruskal groups of
%  (edges of the group joining distinct components BEFORE the group's unions)
%  minus (merges the group performs).  SIX official samples replayed, all
%  hand-verified to 1,0,0,1,0,2.
%  Solved CF 160D "Edges in MST" (2300, L5), live-verified against the CF problem
%  page (title "D. Edges in MST", 2s TL, 256 MB, 2<=n<=1e5,
%  n-1<=m<=min(1e5,n(n-1)/2), 1<=w<=1e6, no loops/multiple edges).  Classify each
%  edge any / at least one / none via per-weight group contraction + Tarjan
%  bridge-finding on the component graph (a bridge = forced = any; a non-bridge
%  usable edge = optional = at least one; endpoints already merged = none).
%  THREE official samples replayed, hand-verified.
%  Extra exports: prob:c17-edges  fig:c17-edges  lst:c17-edges  lem:c17-bridge
% =====================================================================
\section{Which Edges Can Be in an MST?}\label{sec:c17-classify}

\begin{previouslybox}[Ch.31.1 built one MST; now we interrogate all of them]
\(\secref{c17-kruskal}\) returned \emph{a} minimum spanning tree.  Questions
like "must this road be built?", "could this edge ever help?", and "is the MST
unique?" need the full picture: which edges appear in \emph{some} or \emph{every}
MST.  Everything is decided by the same cut and cycle properties, applied inside
each equal-weight group.
\end{previouslybox}

\begin{topicintro}
Sort the edges and slice them into groups of equal weight -- Kruskal's order
never interleaves groups.  Within a group, the components formed by all
\emph{lighter} edges are fixed; the group's job is to merge some of them.  Two
classical lemmas give the trichotomy: the \term{intersection property}
(\(\thmref{c17-intersect}\)) says an edge is in some MST exactly when it is a
minimum edge across the cut it induces; the \term{cycle property}
(\(\thmref{c17-cycle}\)) says an edge that is strictly heaviest on a cycle is in
no MST.  Putting them together splits every edge into forced (in every MST),
optional (in some), or forbidden (in none).

The payoff is a clean count: inside a weight group, let \(\mathit{usable}\) be
the edges joining two \emph{distinct} lighter-components and \(\mathit{merged}\)
the merges Kruskal actually performs there.  Then \(\mathit{usable}-\mathit{merged}\)
is the number of edges that tie for a slot, and the number of MSTs is
\(2^{\text{(free choices)}}\)-style -- we only need the excess count for solved
CF 1108F (\(\probref{c17-unify}\)), which asks the minimum number of \(+1\)
increments needed to make the MST unique without changing its weight.

\textbf{Prerequisites.} \(\secref{c17-kruskal}\) -- Kruskal and its DSU;
\(\thmref{c17-cut}\) and \(\lemref{c17-safe}\).  This section re-runs Kruskal
group by group, so \(\coderef{c17-kruskal}\) is the base being modified.

\textbf{After this section you can}
\begin{itemize}[nosep]
  \item classify every edge as forced / optional / forbidden via the two properties; \hfill\lv{4}
  \item count, per weight group, the edges that tie for an MST slot; \hfill\lv{5}
  \item solve CF 1108F: minimum increments to force a unique MST. \hfill\lv{5}
\end{itemize}
\end{topicintro}

% ---------------------------------------------------------------
\subsection{The two properties and the trichotomy}
% ---------------------------------------------------------------

\begin{definition}{Forced, optional, forbidden}{c17-forced}
An edge is \term{forced} if it lies in every MST, \term{optional} if it lies in
some but not all MSTs, and \term{forbidden} if it lies in none.
\end{definition}

\begin{theorem}{Intersection property}{c17-intersect}
An edge \(e\) belongs to \textbf{some} MST if and only if \(e\) is a minimum-weight
edge crossing the cut defined by removing from the graph every edge of weight
strictly greater than \(w(e)\) and looking at the two components \(e\) would join.
In particular, if \(e\) is the \emph{unique} minimum across its cut it is forced.
\end{theorem}

\begin{theorem}{Cycle property (global form)}{c17-cycle}
An edge is forbidden if and only if it is the strictly heaviest edge on some
cycle.  Equivalently, \(e\) is in no MST iff every path connecting its endpoints
using edges of weight \(\le w(e)\) already exists without \(e\).
\end{theorem}

\begin{lemma}{Equal-weight group excess}{c17-group}
Process the edges by nondecreasing weight in groups.  Just before a group acts,
the DSU reflects all strictly-lighter edges.  In the group, let \(\mathit{usable}\)
count the edges whose two endpoints lie in different components \emph{at that
moment}, and \(\mathit{merged}\) count the successful unions the group performs.
Then \(\mathit{usable}-\mathit{merged}\) edges in this group are \term{optional}
(each is a tie that could replace another in the group), and the excess is added
to the running "non-uniqueness" count.  No edge in the group is forced if it is
part of a same-weight cycle.
\end{lemma}

\begin{lstlisting}[caption={Per-group Kruskal that returns the total optional-excess sum. Edges are pre-sorted; each maximal equal-weight block is scanned twice: once to count usable edges against the pre-group DSU, once to perform the merges}, label={lst:c17-group}]
#include <bits/stdc++.h>
using namespace std;
struct Edge { int u, v; long long w; };
int par[200005];
int findSet(int x){ return par[x]==x ? x : par[x]=findSet(par[x]); }
long long optionalExcess(int n, vector<Edge>& edges){   // edges sorted by w
  for(int i=1;i<=n;i++) par[i]=i;
  long long excess=0; int m=(int)edges.size();
  for(int i=0; i<m; ){
    int j=i; while(j<m && edges[j].w==edges[i].w) j++;  // block [i,j) same weight
    long long usable=0, merged=0;
    for(int k=i; k<j; k++){                              // count BEFORE unioning
      if(findSet(edges[k].u)!=findSet(edges[k].v)) usable++;
    }
    for(int k=i; k<j; k++){                              // now perform merges
      int a=findSet(edges[k].u), b=findSet(edges[k].v);
      if(a!=b){ par[a]=b; merged++; }
    }
    excess += (usable - merged);
    i=j;
  }
  return excess;
}
\end{lstlisting}

% ---------------------------------------------------------------
\subsection{Solved: CF 1108F -- MST Unification (2100, L5)}
% ---------------------------------------------------------------

\solvedproblem{c17-unify}{MST Unification}{CF 1108F}{https://codeforces.com/problemset/problem/1108/F}{CF 2100 L5}{equal-weight Kruskal groups, cycle / intersection properties, DSU, MST uniqueness counting}

\begin{problemstatement}[CF 1108F, verified against the live CF problem page]
A connected weighted graph on \(n\) vertices and \(m\) edges is given (each edge
\(u,v,w\)).  You may repeatedly increase the weight of any single edge by \(1\);
each increment costs one operation.  Find the minimum number of operations needed
so that the graph has a \textbf{unique} minimum spanning tree, \emph{without}
changing the total weight of the MST (you are only allowed to break ties, not to
make the tree more expensive).  Constraints: \(1\le n,m\le 2\cdot 10^5\),
\(1\le w\le 10^9\), graph is connected.  Time limit 3 seconds, memory 256 MB.
There are six sample tests with answers \(1,0,0,1,0,2\): an \(8\)-vertex graph
answer \(1\); a \(4\)-vertex tree answer \(0\); a \(3\)-cycle of weights
\(1,2,3\) answer \(0\); a \(3\)-cycle of weights \(1,3,3\) answer \(1\); a
single vertex \(n=1\) answer \(0\); and a \(5\)-vertex graph answer \(2\).
\end{problemstatement}

\begin{figure}[H]
\centering
\begin{tikzpicture}[diagram, scale=0.9, transform shape,
   tnode/.style={circle, draw=primaryblue!70, fill=primaryblue!8, line width=0.7pt,
                 minimum size=6.5mm, font=\scriptsize\bfseries, inner sep=1pt}]
  \node[tnode] (a) at (0,0) {1};
  \node[tnode] (b) at (2.6,0) {2};
  \node[tnode] (c) at (1.3,-2) {3};
  \draw[successgreen!80!black, line width=1.6pt] (a)--node[fill=white,font=\scriptsize,inner sep=1pt] {1} (b);
  \draw[warningred!80, line width=1.3pt] (b)--node[fill=white,font=\scriptsize,inner sep=1pt] {3} (c);
  \draw[difforange, line width=1.3pt, dashed] (a)--node[fill=white,font=\scriptsize,inner sep=1pt] {3} (c);
  \node[annot, font=\scriptsize, text width=6cm, align=center] at (1.3,-3)
    {sample 4: two weight-3 edges tie; nudge one to 4 (red -> dashed) for uniqueness};
\end{tikzpicture}
\caption{Sample 4, the \(3\)-cycle \(1\!-\!2{:}1,\ 2\!-\!3{:}3,\ 1\!-\!3{:}3\).  The
MST must take the weight-\(1\) edge and exactly one weight-\(3\) edge; the two
weight-\(3\) edges tie, so the MST is not unique.  Bumping one of them to \(4\)
(costs one increment) makes the \(1\!-\!3{:}3\) edge forced and the answer
\(\mathbf{1}\).}
\label{fig:c17-unify}
\end{figure}

\begin{attemptgate}[Attempt it yourself first -- \lv{5}, give yourself 45 minutes]
Uniqueness fails only where two edges of \emph{equal} weight could occupy the same
MST slot.  Rather than trying every increment, what does a single Kruskal run --
grouped by weight -- tell you about how many such ties exist?
\peek{Non-uniqueness comes entirely from equal-weight groups.  In each group,
Kruskal picks a spanning forest among the components formed by lighter edges.
Whenever more edges in the group could have been picked than the number of merges
needed, there is a tie. (\(\secref{c17-classify}\), \(\lemref{c17-group}\).)}
\nudge{Sort edges, walk them in equal-weight blocks.  For a block, first count
\(\mathit{usable}\): edges whose endpoints sit in different DSU components
\emph{before} any union of this block.  Then run the unions and count
\(\mathit{merged}\).  Each of the \(\mathit{usable}-\mathit{merged}\) surplus edges
sits on an equal-weight cycle and must be lifted by one to break the tie, and one
lift fixes exactly one such edge.}
\fullhint{Answer \(=\sum_{\text{groups}}(\mathit{usable}-\mathit{merged})\), computed
by \(\coderef{c17-group}\).  Why one increment per surplus edge suffices and is
necessary: a surplus edge is optional (lies in some MST), and lifting it by \(1\)
removes it from every minimum cut it ties on, turning it forbidden, exactly
resolving one degree of freedom; it never touches the MST weight because a
different same-weight edge still fills the slot. Full build and six-sample replay
below.}
\end{attemptgate}

\solstep{Understanding}
The MST weight is fixed; we are only allowed to re-shuffle \emph{ties} into a
strict order by adding \(1\) to chosen edges.  So we must not perturb lighter
edges' decisions -- only break equal-weight ambiguity.  Grouping Kruskal by
weight isolates exactly that ambiguity: inside a group, the number of edges that
\emph{could} legally merge the lighter-components exceeds the number that
\emph{must} merge them by precisely the count of optional edges, and each needs
one bump.

\solstep{Key Observations}
\begin{observation}
An increment that preserves the MST weight can only raise an edge that is
\emph{not} forced and is currently tied at its weight -- i.e.\ an optional edge
from \(\lemref{c17-group}\).  Raising a lighter or forced edge changes the tree
cost, which is forbidden.
\end{observation}
\begin{observation}
Within one equal-weight group, \(\mathit{usable}\) edges connect distinct
lighter-components, but only \(\mathit{merged}\) of them can be in any spanning
forest of the contracted graph.  The \(\mathit{usable}-\mathit{merged}\) extras
are the interchangeable ones; each extra is one source of non-uniqueness.
\end{observation}
\begin{observation}
A single \(+1\) on one surplus edge removes it from the tie (it becomes strictly
heavier than its group) and thereby makes the remaining choice unique for that
slot.  Different surplus edges need different bumps, so the operations add up:
the minimum total is the sum of per-group excesses.
\end{observation}

\solstep{Think It With Me}
\begin{thinkbox}
\deadend{Recompute "is the MST unique?" after every candidate increment}{MST
uniqueness per check is \(\bigO{m\log m}\); searching increments makes it
quadratic at best -- hopeless at \(m=2\cdot 10^5\).}
\deadend{Count pairs of equal-weight edges}{Equal weight alone is not the
problem: two equal edges that connect four distinct lighter-components are both
forced (no tie).  Sample 1 has many equal weights yet answer \(1\), so raw
equal-edge counts overshoot.}
\aha{Run Kruskal in weight groups and, per group, subtract merges from usable
edges.  The surplus \(\mathit{usable}-\mathit{merged}\) is exactly the number of
optional edges -- and the number of required increments.}
\end{thinkbox}

\solstep{The Approach Ladder}
\begin{center}\renewcommand{\arraystretch}{1.25}\setlength{\tabcolsep}{3pt}
\begin{tabular}{@{}c p{0.28\linewidth} p{0.30\linewidth} l c@{}}
\toprule
\textbf{\#} & \textbf{Approach} & \textbf{Idea} & \textbf{Time} & \textbf{Verdict}\\
\midrule
1 & \lv{3}\ enumerate MSTs & list all MSTs, test uniqueness & exponential & \tle\\
2 & \lv{4}\ count equal-weight pairs & ignore component structure & \(\bigO{m}\) & \textbf{WA}\\
3 & \lv{5}\ \textbf{group usable \(-\) merged} & per weight block excess & \(\bigO{m\log m}\) & \ac\\
\bottomrule
\end{tabular}\end{center}
Rung 2 counts equal-weight edges, but two equal edges on four distinct components
are both forced and create no tie -- it overcounts (fails Sample 1).  Rung 3
respects the lighter-component structure.

\solstep{Solution}
Sort edges by weight.  Walk maximal equal-weight blocks with the DSU frozen at the
pre-block state.  For each block, \(\mathit{usable}\) is the number of block edges
whose endpoints are in different components (checked \emph{before} any union in
the block); then perform the unions counting \(\mathit{merged}\).  Accumulate
\(\mathit{usable}-\mathit{merged}\).  \(\coderef{c17-group}\) is the whole
algorithm; the answer fits in \(64\)-bit but is at most \(m\), so \texttt{int}
would also do for the sum.

\solstep{Implementation}
\begin{lstlisting}[caption={CF 1108F. Sort by weight; per equal-weight block count usable (against the pre-block DSU) then merged (after unioning), and add the difference. Exactly the two-scan loop of optionalExcess}, label={lst:c17-unify}]
#include <bits/stdc++.h>
using namespace std;
struct Edge { int u, v, w; };
int par[200005];
int findSet(int x){ return par[x]==x ? x : par[x]=findSet(par[x]); }
int main(){
  int n, m; scanf("%d %d", &n, &m);
  vector<Edge> edges(m);
  for(int i=0;i<m;i++) scanf("%d %d %d", &edges[i].u, &edges[i].v, &edges[i].w);
  sort(edges.begin(), edges.end(), [](const Edge& a, const Edge& b){ return a.w<b.w; });
  for(int i=1;i<=n;i++) par[i]=i;
  long long ans=0;
  for(int i=0;i<m;){
    int j=i; while(j<m && edges[j].w==edges[i].w) j++;   // block [i,j)
    long long usable=0, merged=0;
    for(int k=i;k<j;k++)
      if(findSet(edges[k].u)!=findSet(edges[k].v)) usable++;   // before unions
    for(int k=i;k<j;k++){
      int a=findSet(edges[k].u), b=findSet(edges[k].v);
      if(a!=b){ par[a]=b; merged++; }                          // do unions
    }
    ans += usable - merged;
    i=j;
  }
  printf("%lld\n", ans);
  return 0;
}
\end{lstlisting}

\solstep{Trace on the Official Samples}
\begin{tracebox}[Sample 4: \, \(3\)-cycle \(1\!-\!2{:}1,\ 2\!-\!3{:}3,\ 1\!-\!3{:}3\) \(\to 1\)]
Group \(w{=}1\): edge \(1\!-\!2\).  usable \(1\) (distinct), merged \(1\) -- excess
\(0\).  Group \(w{=}3\): before it, comps \(\{1,2\},\{3\}\).  Edges \(2\!-\!3\) and
\(1\!-\!3\) both join \(\{1,2\}\) to \(\{3\}\): usable \(=2\).  Unions: first merges
(\(1\)), second now same component (skip): merged \(=1\).  Excess \(2-1=1\).  Total
\(\mathbf{1}\).
\end{tracebox}
\begin{tracebox}[Sample 3: \, \(3\)-cycle \(1\!-\!2{:}1,\ 2\!-\!3{:}2,\ 1\!-\!3{:}3\) \(\to 0\)]
Three groups \(w{=}1,2,3\), each a single edge.  Every group: usable \(=\) merged
\(=1\) (or \(0\) for the last, which closes a cycle: usable \(0\)).  Excess \(0\)
each.  Total \(\mathbf{0}\) -- distinct weights already force a unique MST.
\end{tracebox}
\begin{tracebox}[Sample 6: \, \(n=5\): \(1\!-\!2{:}2,2\!-\!3{:}1,4\!-\!5{:}3,2\!-\!4{:}2,1\!-\!4{:}2,1\!-\!5{:}3\) \(\to 2\)]
\(w{=}1\): \(2\!-\!3\) usable \(1\) merged \(1\) excess \(0\).  \(w{=}2\): block
\(\{1\!-\!2,2\!-\!4,1\!-\!4\}\) over comps \(\{1\},\{2,3\},\{4\},\{5\}\): all three
endpoints distinct at check time \(=\) usable \(3\); unions merge into one
component with \(2\) successful merges, third is a same-weight cycle edge \(=\)
merged \(2\); excess \(1\).  \(w{=}3\): block \(\{4\!-\!5,1\!-\!5\}\) over
\(\{1,2,3,4\},\{5\}\): usable \(2\), merged \(1\), excess \(1\).  Total
\(0+1+1=\mathbf{2}\).
\end{tracebox}
The remaining samples -- Sample 1 (an \(8\)-vertex graph whose only surplus sits in
the weight-\(3\) block) gives \(1\); Sample 2 (a tree) and Sample 5 (\(n=1\)) have
no cycles, so every group merges everything it can: excess \(0\).  All six
answers \(1,0,0,1,0,2\) are reproduced.

\begin{figure}[H]
\centering
\begin{tikzpicture}[diagram, scale=0.9, transform shape,
   comp/.style={circle, draw=primaryblue!70, fill=primaryblue!8, line width=0.8pt,
               minimum size=7mm, font=\scriptsize\bfseries, inner sep=1pt},
   grp/.style={rectangle, rounded corners=2pt, draw=difforange!70, fill=difforange!10,
               line width=0.9pt, inner sep=2pt, font=\scriptsize}]
  \node[comp] (A) at (0,0) {\(\{1\}\)};
  \node[comp] (B) at (3,0) {\(\{2,3\}\)};
  \node[comp] (C) at (1.5,2) {\(\{4\}\)};
  \draw[difforange, line width=1.3pt] (A)--node[grp] {1-2} (B);
  \draw[difforange, line width=1.3pt] (B)--node[grp] {2-4} (C);
  \draw[difforange, line width=1.3pt, dashed] (A)--node[grp] {1-4} (C);
  \node[annot, font=\scriptsize, text width=8cm, align=center] at (1.5,-1.1)
    {weight-2 block: usable 3, merged 2, excess 1 (the dashed edge is optional)};
\end{tikzpicture}
\caption{Sample 6's weight-\(2\) group.  The three lighter-components
\(\{1\},\{2,3\},\{4\}\) need only \(2\) merges, but \(3\) edges offer to connect
them -- the triangle.  One edge is surplus (dashed) and must be lifted by \(1\).}
\label{fig:c17-groups}
\end{figure}

\solstep{Complexity}
\begin{complexitybox}
Sorting is \(\bigO{m\log m}\); the two-scan group loop is \(\bigO{m\,\alpha(n)}\)
with path-compressed DSU (\chref{c28}).  Total \(\bigO{m\log m}\) time,
\(\bigO{n+m}\) space -- comfortable for \(n,m\le 2\cdot 10^5\) in 3 seconds.  The
counters are conceptually \(\bigO{m}\), but \texttt{long long} is used for the
answer to be safe.
\end{complexitybox}

\begin{hardeningbox}
\begin{itemize}[nosep,leftmargin=1.4em]
\item \textbf{Two scans are mandatory.}  Counting \(\mathit{usable}\) must read the
      DSU \emph{before} any union in the block; merging in the same loop (single
      scan) makes later edges of the block see already-merged components and
      undercounts usable \(--\) a classic WA.
\item \textbf{Block boundary.}  Group strictly by equal weight; a graph with
      \(\bigO{m}\) distinct weights degenerates to single-edge blocks where excess
      is always \(0\), matching "distinct weights \(\Rightarrow\) unique MST".
\item \textbf{Disconnected check.}  The problem guarantees connectedness; if you
      reuse \(\coderef{c17-group}\) elsewhere, verify \(\mathit{merged}\) totals
      \(n-1\).
\item \textbf{Overflow.}  The answer is at most \(m\le 2\cdot 10^5\); safe in
      \texttt{int}, but \texttt{long long} avoids surprises if constraints grow.
\end{itemize}
\end{hardeningbox}

% ---------------------------------------------------------------
\subsection{Solved: CF 160D -- Edges in MST (2300, L5)}
% ---------------------------------------------------------------

\solvedproblem{c17-edges}{Edges in MST}{CF 160D}{https://codeforces.com/problemset/problem/160/D}{CF 2300 L5}{equal-weight group contraction, Tarjan bridge-finding, cut / cycle properties, MST edge classification}

\begin{problemstatement}[CF 160D, verified against the live CF problem page]
A connected weighted undirected graph on \(n\) vertices and \(m\) edges (no loops,
no multiple edges) is given, each edge \(a,b,w\).  For \textbf{each} edge, decide
whether it is in \emph{every} MST (print \texttt{any}), in \emph{at least one} MST
(print \texttt{at least one}), or in \emph{no} MST (print \texttt{none}), reporting
in input order.  Constraints: \(2\le n\le 10^5\), \(n-1\le m\le\min(10^5,
n(n-1)/2)\), \(1\le w\le 10^6\).  Time limit 2 seconds, memory 256 MB.  Samples:
\(4\ 5\) with edges \((1,2,101),(1,3,100),(2,3,2),(2,4,2),(3,4,1)\) \(\to\)
\texttt{none,any,at least one,at least one,any}; \(3\ 3\) edges
\((1,2,1),(2,3,1),(1,3,2)\) \(\to\) \texttt{any,any,none}; \(3\ 3\) edges
\((1,2,1),(2,3,1),(1,3,1)\) \(\to\) \texttt{at least one,at least one,at least one}.
\end{problemstatement}

\begin{lemma}{Bridge \(\iff\) forced inside a group}{c17-bridge}
Fix an equal-weight group and contract every strictly-lighter edge, so the group's
usable edges become the edges \(H\) of a multigraph on the current components.  A
usable edge lies in \emph{every} MST of \(G\) (is \texttt{any}) \(\iff\) it is a
\term{bridge} of \(H\): deleting it disconnects \(H\), so it is the unique cheapest
crossing of the cut its removal defines (\(\thmref{c17-intersect}\)).  A usable
non-bridge sits on a cycle of \(H\) and is therefore \texttt{at least one}.
\end{lemma}

\begin{figure}[H]
\centering
\begin{tikzpicture}[diagram, scale=0.9, transform shape,
   comp/.style={circle, draw=primaryblue!70, fill=primaryblue!8, line width=0.8pt,
               minimum size=9mm, font=\scriptsize\bfseries, inner sep=1pt},
   iso/.style={circle, draw=gray!70, fill=gray!8, line width=0.8pt, dashed,
               minimum size=9mm, font=\scriptsize\bfseries, inner sep=1pt},
   nb/.style={line width=1.4pt, draw=difforange}]
  \node[iso] (A) at (6,0) {\(\{1\}\)};
  \node[comp] (B) at (0,0.9) {\(\{2\}\)};
  \node[comp] (C) at (0,-0.9) {\(\{3,4\}\)};
  \draw[nb,bend left=16] (B)--node[right,font=\scriptsize] {2-3} (C);
  \draw[nb,bend right=16] (B)--node[left,font=\scriptsize] {2-4} (C);
  \node[annot, font=\scriptsize, text width=5cm, align=center] at (6,-1.3)
    {\(\{1\}\) has no \(w{=}2\) edge -- isolated in \(H\)};
  \node[annot, font=\scriptsize, text width=8cm, align=center] at (0,-2.7)
    {two parallel edges form a \(2\)-cycle: neither is a bridge \(\to\) both at least one};
\end{tikzpicture}
\caption{Sample 1's weight-\(2\) group after contracting the lighter
\(3\!-\!4{:}1\) edge: components \(\{1\},\{2\},\{3,4\}\).  The only usable edges are
\(2\!-\!3\) and \(2\!-\!4\), which are \emph{parallel} in \(H\) (a \(2\)-cycle); a bridge must
lie on no cycle, so neither is a bridge and both are \texttt{at least one}.  Vertex
\(\{1\}\) has no weight-\(2\) edge, so it is isolated.  Contrast the lone \(w{=}1\) edge
\(3\!-\!4\) and the lone \(w{=}100\) edge \(1\!-\!3\), each a single bridge in its group and
therefore \texttt{any}.}\label{fig:c17-edges}
\end{figure}

\begin{attemptgate}[Attempt it yourself first -- \lv{5}, give yourself 45 minutes]
A single Kruskal marks the edges it picked, but several MSTs exist -- "picked once"
is not "in every MST."  You need three answers per edge.  Reduce to one weight
group at a time: after contracting lighter edges, what graph property separates a
forced edge from a merely-optional one?
\peek{Inside a group the lighter components are fixed.  Draw the group's edges on
those contracted components -- a multigraph \(H\).  Forced edges are exactly the
bridges of \(H\); usable non-bridges are optional; edges whose endpoints were
already joined are forbidden. (\(\lemref{c17-bridge}\), \(\thmref{c17-cycle}\).)}
\nudge{Run group-Kruskal.  For each weight block, build \(H\) from the edges whose
endpoints lie in distinct pre-block components, then find \(H\)'s bridges with
Tarjan's \(\mathrm{low}/\mathrm{tin}\) DFS.  A usable edge that is a bridge
\(\to\) \texttt{any}; a usable non-bridge \(\to\) \texttt{at least one}; an edge with
both endpoints already in one component \(\to\) \texttt{none}.  Only union the block
\emph{after} classifying it.}
\fullhint{Because \(H\) is a multigraph, run bridge-finding keyed on the \emph{edge
identity} (skip the exact reverse arc, not the parent vertex) so parallel edges are
correctly seen as non-bridges.  Classify before merging so all block edges see the
same frozen pre-block DSU.  Total \(\bigO{m\log m}\): the per-group Tarjan work
sums to \(\bigO{m}\) over all groups.  Full code and three-sample replay below.}
\end{attemptgate}

\solstep{Understanding}
The trichotomy is decided \emph{inside each equal-weight group} once lighter edges
are contracted.  An edge forbidden by \(\thmref{c17-cycle}\) has its endpoints
already joined before its group.  Among usable edges, \(\lemref{c17-bridge}\) makes
forced-vs-optional a pure graph-connectivity question -- bridges -- on the
contracted group graph.  So the algorithm is group-Kruskal wrapped around one
bridge-finding pass per group.

\solstep{Key Observations}
\begin{observation}
An edge is \texttt{none} iff, at the moment its weight group is about to run,
its endpoints are already in the same DSU component -- it is the strictly heaviest
edge on a cycle of lighter-or-equal edges (\(\thmref{c17-cycle}\)).
\end{observation}
\begin{observation}
Among usable edges of a group, contracting the lighter components turns the group
into a multigraph \(H\).  A bridge of \(H\) is the unique lightest crossing of the
cut it induces, hence \texttt{any}; a non-bridge lies on a cycle of equal-weight
edges, so an MST exists using it and an MST exists avoiding it -- \texttt{at least
one} (\(\lemref{c17-bridge}\)).
\end{observation}
\begin{observation}
Bridge-finding on a \emph{multigraph} must distinguish parallel edges by identity:
skipping only the arc you came along (an edge slot), not every arc to the parent
vertex, otherwise a pair of parallel edges is mis-flagged as two bridges.
\end{observation}

\solstep{Think It With Me}
\begin{thinkbox}
\deadend{Run Kruskal once and print \texttt{any} for taken edges, \texttt{none} for the rest}{Taken-once is neither forced nor even ``in some MST''; Sample 3 has every edge equal weight -- every pair forms an MST, so all are \texttt{at least one}, and single-Kruskal mislabels two of them \texttt{any}.}
\deadend{For each edge, re-run Kruskal forcing and forbidding it}{\(\bigO{m}\) MST runs at \(\bigO{m\log m}\) each is \(\bigO{m^2\log m}\) -- \emph{TLE} at \(m=10^5\).}
\aha{Contract lighter edges per weight group, then a single Tarjan bridge pass on
the group's component graph separates bridges (\texttt{any}) from cycle edges
(\texttt{at least one}); pre-merged endpoints are \texttt{none}.}
\end{thinkbox}

\solstep{The Approach Ladder}
\begin{center}\renewcommand{\arraystretch}{1.25}\setlength{\tabcolsep}{3pt}
\begin{tabular}{@{}c p{0.28\linewidth} p{0.30\linewidth} l c@{}}
\toprule
\textbf{\#} & \textbf{Approach} & \textbf{Idea} & \textbf{Time} & \textbf{Verdict}\\
\midrule
1 & \lv{3}\ one Kruskal & taken \(\to\) any, else none & \(\bigO{m\log m}\) & \textbf{WA}\\
2 & \lv{4}\ force/forbid per edge & rerun MST \(m\) times & \(\bigO{m^2\log m}\) & \tle\\
3 & \lv{5}\ \textbf{group + bridges} & contract, Tarjan \(H\) & \(\bigO{m\log m}\) & \ac\\
\bottomrule
\end{tabular}\end{center}
Rung 1 ignores ties (fails Sample 3); rung 2 is correct but quadratic.  Rung 3 uses
the bridge characterisation to classify a whole group in linear time.

\solstep{Solution}
Sort edges by weight and walk maximal equal-weight blocks with the DSU frozen at
the pre-block state.  For the block, discard edges whose endpoints are already
together (\texttt{none}); renumber the surviving components and add the remaining
(usable) edges as slots of a multigraph \(H\).  Run Tarjan bridge-finding on \(H\): a
usable edge that is a bridge is \texttt{any}, otherwise \texttt{at least one}.  Only
after classifying the block, union \emph{all} its usable edges so the next group
sees the merged components.  Print the label per original edge index.

\solstep{Implementation}
\begin{lstlisting}[caption={CF 160D. Per weight group: contract lighter edges, build the multigraph H of usable edges, Tarjan bridges (skip only the reverse edge-slot, so parallel edges are never false bridges). bridge->any, non-bridge usable->at least one, pre-merged->none. Union the whole block only after classifying}, label={lst:c17-edges}]
#include <bits/stdc++.h>
using namespace std;
struct Edge { int u, v, w, id; };
const int MAXN = 100005;
int dsu[MAXN];
int findSet(int x){ return dsu[x]==x ? x : dsu[x]=findSet(dsu[x]); }
int res[MAXN];                     // 0 none, 1 at least one, 2 any
vector<pair<int,int> > adj[MAXN];  // (compact neighbour, edge-slot) for H
int tinA[MAXN], lowA[MAXN], clk;
vector<int> isBridge;
void dfsBridge(int v, int pe){
  tinA[v]=lowA[v]=++clk;
  for(size_t i=0;i<adj[v].size();i++){
    int to=adj[v][i].first, eid=adj[v][i].second;
    if(eid==pe) continue;                       // skip ONLY the reverse slot
    if(tinA[to]) lowA[v]=min(lowA[v], tinA[to]);
    else { dfsBridge(to,eid); lowA[v]=min(lowA[v],lowA[to]);
           if(lowA[to]>tinA[v]) isBridge[eid]=1; }
  }
}
int main(){
  int n,m; scanf("%d %d",&n,&m);
  vector<Edge> e(m);
  for(int i=0;i<m;i++){ scanf("%d %d %d",&e[i].u,&e[i].v,&e[i].w); e[i].id=i; }
  sort(e.begin(), e.end(), [](const Edge&a,const Edge&b){return a.w<b.w;});
  for(int i=1;i<=n;i++) dsu[i]=i;
  for(int i=0;i<m;){
    int j=i; while(j<m && e[j].w==e[i].w) j++;            // block [i,j)
    unordered_map<int,int> comp; int cn=0;
    vector<pair<int,int> > us;                            // (ca,cb) per slot
    vector<int> usId;                                      // original edge id per slot
    for(int k=i;k<j;k++){
      int a=findSet(e[k].u), b=findSet(e[k].v);
      if(a==b){ res[e[k].id]=0; continue; }               // none
      if(!comp.count(a)) { comp[a]=cn; adj[cn].clear(); cn++; }
      if(!comp.count(b)) { comp[b]=cn; adj[cn].clear(); cn++; }
      us.push_back(make_pair(comp[a], comp[b])); usId.push_back(e[k].id);
    }
    for(int s=0;s<(int)us.size();s++){                    // build H (multi-edge)
      int ca=us[s].first, cb=us[s].second;
      adj[ca].push_back(make_pair(cb,s)); adj[cb].push_back(make_pair(ca,s));
    }
    isBridge.assign(us.size(),0); clk=0;
    for(int t=0;t<cn;t++) tinA[t]=0;
    for(int t=0;t<cn;t++) if(!tinA[t]) dfsBridge(t,-1);    // Tarjan per component
    for(int s=0;s<(int)us.size();s++)
      res[usId[s]] = isBridge[s] ? 2 : 1;                 // bridge->any, else at least one
    i=j;
  }
  for(int id=0;id<m;id++)
    puts(res[id]==2?"any":(res[id]==1?"at least one":"none"));
  return 0;
}
\end{lstlisting}

\solstep{Trace on the Official Samples}
\begin{tracebox}[Sample 1: \, \(4\) vertices, edges \(1\!-\!2{:}101,\ 1\!-\!3{:}100,\ 2\!-\!3{:}2,\ 2\!-\!4{:}2,\ 3\!-\!4{:}1\)]
\(w{=}1\): \(3\!-\!4\) joins \(\{3\},\{4\}\) -- single bridge \(\to\) \texttt{any}; union
\(\{3,4\}\).  \(w{=}2\): comps \(\{1\},\{2\},\{3,4\}\); edges \(2\!-\!3\) (\(\{2\}\!-\!\{3,4\}\))
and \(2\!-\!4\) (\(\{2\}\!-\!\{3,4\}\)) are \emph{parallel} in \(H\) -- a \(2\)-cycle, no
bridge \(\to\) both \texttt{at least one}; union \(\{2,3,4\}\).  \(w{=}100\): \(1\!-\!3\)
joins \(\{1\},\{2,3,4\}\), single bridge \(\to\) \texttt{any}.  \(w{=}101\): \(1\!-\!2\)
endpoints already together \(\to\) \texttt{none}.  Input order:
\texttt{none,any,at least one,at least one,any}. \(\checkmark\)
\end{tracebox}
\begin{tracebox}[Sample 2: \, \(1\!-\!2{:}1,\ 2\!-\!3{:}1,\ 1\!-\!3{:}2\)]
\(w{=}1\): \(H\) on \(\{1\},\{2\},\{3\}\) has edges \(1\!-\!2,\ 2\!-\!3\) -- a path, both
bridges \(\to\) \texttt{any,any}; union all.  \(w{=}2\): \(1\!-\!3\) endpoints together
\(\to\) \texttt{none}.  Output \texttt{any,any,none}. \(\checkmark\)
\end{tracebox}
\begin{tracebox}[Sample 3: \, triangle \(1\!-\!2{:}1,\ 2\!-\!3{:}1,\ 1\!-\!3{:}1\)]
One group \(w{=}1\): \(H\) is the triangle on \(\{1\},\{2\},\{3\}\) -- no bridge, so all
three are \texttt{at least one}.  Output \texttt{at least one, at least one, at least
one}. \(\checkmark\)
\end{tracebox}

\solstep{Complexity}
\begin{complexitybox}
Sorting is \(\bigO{m\log m}\).  Each edge enters exactly one group's \(H\) and is
touched once by the two-scan plus the Tarjan DFS, whose total work over all groups
is \(\bigO{n+m}\).  Hence \(\bigO{m\log m}\) time overall and \(\bigO{n+m}\) space --
comfortable for \(n,m\le 10^5\) in 2 seconds.  The strings are printed verbatim
lowercase; a capital \texttt{Any} is a presentation-error WA.
\end{complexitybox}

\begin{hardeningbox}
\begin{itemize}[nosep,leftmargin=1.4em]
\item \textbf{Classify before unioning} the block: every usable edge must see the
      same frozen pre-block DSU, or the group's \(H\) collapses and forced edges look
      forbidden.
\item \textbf{Multigraph bridges:} skip only the reverse edge-slot, never every arc
      to the parent vertex; parallel edges must come out non-bridges (Sample 1's
      \(w{=}2\) pair).
\item \textbf{Reset per group:} clear the compact adjacency and re-initialise
      \(\mathrm{tin}\)/timestamp for the \(0..cn{-}1\) nodes you use.
\item \textbf{Forbidden first:} endpoints already merged \(\to\) \texttt{none}
      must be excluded from \(H\) before bridge-finding, not classified afterwards.
\end{itemize}
\end{hardeningbox}

\begin{recallbox}[Recall -- what \(\secref{c17-classify}\) gave you]
\selftest{
  \item State the intersection and cycle properties in one line each.
  \item Why can two equal-weight edges impose \emph{no} tie?
  \item Give the per-group formula for the optional-edge excess.
  \item Why does one \(+1\) increment fix exactly one surplus edge, and never the MST weight?
  \item In CF 160D, why is a bridge of the contracted group graph exactly a \texttt{any} edge?
}
\end{recallbox}
